% sections/06-exploitation.tex — AIxCC treated as competition + system-design case (spec §6).
\section{Exploitation and Validation: The AIxCC Case Study}
\label{sec:aixcc}

The Defense Advanced Research Projects Agency's Artificial Intelligence Cyber Challenge
(AIxCC) is the only included setting in our corpus where discovery, proof-of-vulnerability
generation, and patching were integrated, scored, and operated without human intervention --- and it is best
analyzed not as a benchmark but as a deliberately constructed offense$\leftrightarrow$defense
instrument whose rules, incentives, and aftermath document how the field's unit of progress was
administratively redefined (I1).

\subsection{Design and Incentives}
Announced in 2023 with ARPA-H participation and support from Anthropic, Google, Microsoft, and
OpenAI, the challenge funneled 42 accepted semifinal teams down to seven finalists, each
receiving \$2M to harden a Cyber Reasoning System (CRS) for a final round ``set up to mimic a
real world CI/CD pipeline''~\cite{openssfaixcc2026}. Scoring rewarded finding vulnerabilities,
proving them via PoVs, and applying correct patches, with speed and accuracy modifiers;
``human interaction was strictly prohibited'' during scored operation~\cite{tobbuttercup2025}.
Two incentive details matter for interpretation. Patching was worth three times discovery
(6 vs.\ 2 points), making B3 capability the rational investment~\cite{taasc2024}; and accuracy
modifiers penalized incorrect patches quadratically --- Team Atlanta reports its PoV-validated
91.27\% patch accuracy yielding a modifier of $1-(1-0.9127)^4=0.9999$ against Theori's
PoV-free $44.4\%$ yielding $0.9044$~\cite{taafc2025}. The scoring mathematics thus priced
patch-quality adjudication directly into strategy.

\subsection{What the Final Round Measured}
Participant and systematization primaries agree on the scored structure: 48 challenge projects
across 23 open-source repositories in the final round~\cite{tobbuttercup2025}, operated over
what the systematization records as ``143 hours of fully autonomous operation'' by CRSs from
all seven finalist teams~\cite{r7sokaixcc2026}. Crucially, the CRSs also surfaced real,
non-synthetic vulnerabilities --- Team Atlanta reports six C/C++ and twelve Java bugs found
collectively during the final, including a SQLite zero-day first seen at
semifinals~\cite{taafc2025} --- and Trail of Bits demonstrated a PoV for a vulnerability that
had \emph{not} been inserted into any challenge~\cite{tobbuttercup2025}. Real-bug discovery did
not score --- Team Atlanta footnotes that their autonomous find-and-patch of a real SQLite bug
during semifinals earned nothing competitively~\cite{taasc2024} --- yet it is precisely this
unscored spillover that connects the competition to the production ecosystem. DARPA's official
closeout reports the remaining aggregates: competitors' systems discovered 54 of 63 synthetic
vulnerabilities (86\%) and patched 43 (68\%), alongside 11 patches for the real
vulnerabilities~\cite{darparesults2025}; DARPA's scoring guide separately confirms both the
three-to-one patching incentive and the \$8.5M final prize pool~\cite{darpascoring2025}. The
closeout also carries an editor's note correcting its own earlier count of 70 synthetic
vulnerabilities to 63 --- an instance of administrative record instability we separate from
measurement validity throughout (Sec.~\ref{sec:validity}).

\subsection{CRS Design Space}
The finalist architectures span the design axes our framework distinguishes.
\emph{Atlantis} (Team Atlanta) ran N-version programming across orthogonal sub-CRSs (C, Java,
multilanguage, patching, SARIF), ensemble fuzzing and concolic execution alongside LLM agents,
and an explicit integration taxonomy --- LLM-augmented, LLM-opinionated, LLM-driven ---
with smaller models often outperforming larger or reasoning models for patch tasks; its
validation oracle re-runs PoVs against patched builds while acknowledging that MTE/PAC
recompilation or broad \texttt{catch(Exception)} wrappers can suppress PoVs without semantic
repair~\cite{taafc2025}. \emph{Buttercup} (Trail of Bits) took the opposite bet:
deterministic, expertise-decomposed workflows with LLMs only where classical tools fail,
using exclusively non-reasoning models --- achieving the field's best cost efficiency at
\$181 per point (versus Atlanta's \$263) on 28 vulnerabilities and 19 patches across 20 CWE
classes with ${>}90\%$ accuracy~\cite{tobbuttercup2025,r7sokaixcc2026}. \emph{Theori}
demonstrated purely static, PoV-free patching, trading accuracy for robustness to absent
proofs~\cite{taafc2025}. Team-level compute costs are independently consistent across our two
sources: totals of \$103.3K (Atlanta), \$39.6K (Trail of Bits), and \$31.8K (Theori) appear in
both the systematization analysis and the participant account~\cite{r7sokaixcc2026,tobbuttercup2025}.

\subsection{Brittleness as Environment Evidence}
The most cited lesson inside the winning team is a string-matching heuristic: Atlantis skipped
patch generation for paths containing ``fuzz,'' and organizers' ``ossfuzz'' directory prefixing
nearly disabled the system hours before the deadline --- ``here we were, building an autonomous
CRS powered by state-of-the-art AI, nearly defeated by a string matching bug''~\cite{taafc2025}.
For our validity taxonomy this is textbook V4: benchmark-specific heuristics functioned as
load-bearing components, invisible until the environment shifted.

\subsection{Interpretive Discipline}
Three boundaries govern what AIxCC evidence can support. First, E-class: competition forks with
inserted bugs are environment E1; substantial autonomy under prohibition of human interaction
still fails A3's production-environment condition. Second, outcome definitions: ``discovered''
means a scored synthetic-vulnerability PoV, ``patched'' means oracle-passing --- the oracle-gaming
taxonomy above bounds what ``correctly patched'' certifies. Third, record reliability: figures
here cite corrected primaries (DARPA closeout, participant blogs cross-checked against the
systematization), never press intermediaries. With those bounds, AIxCC demonstrates that
integrated $\{B1{+}B2{+}B3\}$ operation at A2 is achievable within one engineering generation
--- and its aftermath, examined in Sec.~\ref{sec:coevolution}, shows the loop continuing
through funded redeployment rather than ending at a scoreboard.
